\chapter{Будущее: память как динамика (резервуар, ассоциативные аттракторы, кванты)}
\label{ch:dynamical}

\begin{pitfall}[Вердикт вперёд: это смена субстрата, а не смена топологии]
Приложения~\ref{ch:hypergraphs} и~\ref{ch:cellcomplex} спрашивали «какой \emph{формы} должен
быть граф». Это приложение спрашивает радикальнее: \emph{а нужен ли символический граф
вообще}, или память стоит держать как \emph{динамику} — состояние большой рекуррентной
системы, где факт не хранится объектом, а «протоптан» в связях. Идея красивая, глубоко
укоренённая в нейронауке, и \textbf{мы её сейчас не строим}. Но не потому, что она наивна —
наоборот, разберём её всерьёз. Причина в том, что для \emph{компаньона} с требованиями
редактируемости, забывания-по-команде, провенанса и поддержки истинности распределённая
память в весах ломает ровно те свойства, ради которых мы и завели типизированный граф. Ниже
— честный разбор: что здесь ценно (и что мы, оказывается, уже делаем), что дорого, и как
оставить дверь открытой, не переписывая мозг.
\end{pitfall}

\section{Постановка: память в весах, воспоминание как траектория}

Наш субстрат (гл.~\ref{ch:substrate}) — типизированные узлы и явные рёбра: у «узла Helena»
есть содержание, у ребра — тип и вес. Альтернатива, которую предлагает эта ветка, меняет
само определение:
\[
\text{память} = (G, W, \theta),\qquad
\text{воспоминание} = \text{траектория в пространстве состояний},\qquad
\text{обучение} = \Delta W .
\]
Здесь $G$ — топология (какие пути вообще есть), $W$ — веса связей (насколько они
протоптаны), а $\theta$ — параметры узлов-вычислителей. Узел \texttt{\#48372} не значит
«Лиссабон»; смысл распределён по тысячам узлов и существует только как \emph{паттерн
активации} $+$ \emph{связность} $+$ \emph{траектория}. Запрос не «ищет chunk», а
\emph{возмущает} систему: бросаем заряд, сеть эволюционирует, из неё выныривает активный
подграф, и уже его вербализует LLM.

\begin{intuition}[Бассейн, камни, датчики]
Классический образ резервуарных вычислений: бассейн воды, в который бросают камни (входной
поток). Каждый камень поднимает волны; следующий приходит, пока прежние не улеглись, — и
поверхность в момент $t$ несёт отпечаток нескольких последних бросков. Датчики по краям
считывают состояние, и простой линейный слой учится по ним отвечать. Сам бассейн обучать не
надо: его физика уже превращает историю входов в богатое состояние. Reservoir computing —
это \emph{случайное нелинейное динамическое отображение признаков для последовательностей},
аналог kernel trick, но во времени.
\end{intuition}

Важная оговорка сразу: \emph{классический} резервуар (Echo State Network, Liquid State
Machine) держит $W_{\text{res}}$ \emph{замороженным} и обучает только считывающий слой. Он
обладает \emph{затухающей} памятью (fading memory) на конечный горизонт — прекрасно для
временных рядов, но не для «узнала в 2026-м, вспомнила в 2031-м». Для памяти длиной в годы
нужен \emph{пластичный} резервуар ($W$ медленно меняется опытом) плюс механизм устойчивого
хранения из мира \emph{ассоциативных аттракторов} (Hopfield/CA3): воспоминание — это
\emph{бассейн притяжения}, частичный cue скатывается в него (pattern completion). То есть это
не «reservoir computing» в чистом виде, а гибрид на пересечении резервуара, ассоциативной
памяти, аттракторных сетей и пластичной динамики.

\section{Математическая ось: один лапласиан — три динамики}

Самое элегантное в этой ветке — то, что \emph{одна и та же} топология допускает три
принципиально разные динамики. Для графа с лапласианом $L=D-A$:
\begin{align}
\text{диффузия:}\quad & \dot x = -Lx & \Rightarrow\quad & x(t)=e^{-tL}x(0),\\
\text{волна:}\quad & \ddot x = -Lx & \Rightarrow\quad & x(t)=\cos(t\sqrt{L})\,x(0),\\
\text{квантовое блуждание:}\quad & i\dot\psi = L\psi & \Rightarrow\quad & \psi(t)=e^{-itL}\psi(0).
\end{align}
Один $L$ порождает $e^{-tL}$ (рассеивающую память), $\cos(t\sqrt L)$ (волновую, с
интерференцией и без затухания) и $e^{-itL}$ (унитарную квантовую). Это не словесная
аналогия, а линейная алгебра: «бросить заряд в узел и смотреть, куда потечёт» имеет буквальный
волновой и квантовый аналоги на той же сети.

\begin{intuition}[Мы уже живём в углу $e^{-tL}$]
Полезно осознать: наш ретрив (гл.~\ref{ch:walker}) — это \emph{и есть} диффузионный угол
этой оси. `diffuse()` считает усечённое $e^{-tL}$-подобное распространение из seed-якорей;
«вспышки, слабеющие к окраинам» — это спад $e^{-tL}$. То есть значительная часть красоты
этого приложения у нас уже реализована, просто в скромной, редактируемой форме. Волновой и
квантовый углы — это \emph{другие операторы над той же топологией}, и их честно можно
сравнить на одном графе (см.\ открытые вопросы в конце).
\end{intuition}

\section{Нейронаучная опора (принцип, а не крысиный мозг)}

Ветка опирается на реальные результаты, и стоит назвать их аккуратно. Обучение оставляет
устойчивые следы в самой сети (engram-клетки: клетки, участвовавшие в формировании
воспоминания, можно искусственно возбудить и вызвать его; хранение связывают с синаптическим
усилением engram-ансамблей). Гиппокамп работает не как «жёсткий диск», а как система
связывания/индексации и реактивации распределённых корковых представлений (hippocampal
indexing theory), причём со временем память дрейфует от гиппокампа к коре (systems
consolidation, replay). Complementary Learning Systems (CLS) прямо противопоставляет
\emph{быстро} обучающийся гиппокамп \emph{медленно} обучающейся коре — и это вычислительно
мотивированная защита от катастрофического забывания. Аттракторная память (Hopfield; CA3 как
autoassociative network) хранит паттерн как бассейн притяжения. Мы заимствуем \emph{принцип}
(быстрый индекс $+$ медленное распределённое хранилище $+$ восстановление по частичному cue),
не пытаясь эмулировать биологию до деталей.

\section{Предлагаемая гибридная архитектура}

Синтез, к которому сходится обсуждение, — не «резервуар вместо графа», а шестислойный объект
плюс архив сбоку:
\begin{enumerate}\itemsep2pt
\item \textbf{Event encoder.} Приходит не «факт», а \emph{событие} целиком — контент, говорящий,
тон, контекст, время, настроение — и кодируется в вектор $x_t$. Источник знания становится
\emph{частью паттерна}, а не строкой БД (значит, память сможет и \emph{ошибиться} в источнике
— по-человечески).
\item \textbf{Hippocampal pattern separator / index.} Из $x_t$ строится разреженный код
$h_t=\mathrm{kWTA}(Px_t)$ (активны немногие units из многих); похожие события разводятся в
непохожие коды (pattern separation). Это «умный стартер»: куда бросить заряд.
\item \textbf{Fast associative memory.} Быстрая Hebbian-запись «эти вещи произошли вместе»
($H\!\leftarrow\!H+\eta_h h_t h_t^\top$), связывающая индекс с распределённым состоянием.
\item \textbf{Slow plastic cortical reservoir.} Большая разреженная рекуррентная сеть
\[
r(t{+}1)=(1-\alpha)\,r(t)+\alpha\,\sigma\!\big(Wr(t)+Ux(t)+Bh(t)\big),
\]
где $W$ \emph{медленно} пластичен ($\eta_c\ll\eta_h$). Быстрый гиппокамп $+$ медленная кора —
это и есть CLS.
\item \textbf{Continuous activation field.} Живое поле $r(t)$ — «о чём Рин сейчас думает».
\item \textbf{Latent-to-LLM interface.} Мост от $r(t)$ к языковой модели.
\end{enumerate}
Плюс \textbf{replay/consolidation} в простое (ночью гиппокамп многократно проигрывает след,
медленно перестраивая $W$ — «переваривание дня» не через \texttt{summarize()}, а через
проигрывание), и \textbf{граф-архив} как «чёрный ящик самолёта» сбоку.

\begin{principle}[Граф не выбрасывается — он становится архивом и учителем]
Ключевой примирительный ход: наш явный граф перестаёт быть \emph{моделью} человеческой памяти
и становится (а) \emph{эпизодическим архивом} (что реально произошло, когда, откуда знаем) и
(б) \emph{учителем} для пластичной памяти (материал для replay-консолидации). Тогда
специализация чистая: \emph{граф} — «что действительно было записано», \emph{резервуар} — «что
Рин из этого сейчас помнит» (неточно, ассоциативно, с интерференцией). Это не дублирование, а
разделение ролей.
\end{principle}

\section{Настроение как нейромодулятор и «пирожки Алисы»}

Отдельная и, на мой взгляд, самая \emph{сразу применимая} идея ветки: mood перестаёт быть
вектором, который кладут в промпт, и становится \emph{нейромодулятором} — он меняет
\emph{геометрию} recall:
\[
W_{\text{effective}} = \gamma(m)\,W,\qquad
A_{\text{effective}} = A\odot M(\mathrm{mood}).
\]
Тот же cue при разном настроении уходит по разным траекториям («помнишь Тбилиси?» $\to$ переезд
$\to$ Ада; или $\to$ ОИЯИ $\to$ смена жизни). Инерция настроения тут — не недостаток, а
\emph{нужное} свойство: медленный low-pass-контроллер задаёт устойчивый ландшафт для быстрых
процессов и не «дёргает руль» на каждый токен.

\begin{intuition}[Гомеостаз как гормоны: заманчиво, но чинит болезнь, которую само же создаёт]
Предложение «мониторить энергию поля $E(t)=\frac1N\sum r_i^2$ и впрыскивать седатив (поднять
утечку $\lambda$, снизить крутизну $\beta$) при перевозбуждении или кофеин при апатии» — это
настоящая идея из neuromodulated networks, и она красива. Но будем честны: это \emph{лекарство
от нестабильности, которую вносит сам пластичный резервуар}. На нашем текущем субстрате
диффузия ограничена по построению (decay $+$ floor $+$ max\_steps $+$ max-CRDT активация,
гл.~\ref{ch:walker}) — «эпилептического шторма» просто негде взяться, и гормональный
контроллер не нужен. Более того, он добавляет ровно то, чего мы весь проект избегали:
свободные параметры ($E_{\min},E_{\max}$, коэффициенты впрыска), которые нечем калибровать без
downstream-сигнала. Вывод сбалансированный: \emph{mood-как-нейромодулятор} — берём хоть сейчас
(он и так входит в правило перехода через $\mathrm{mc}(B\mid\mood)$); \emph{гормональный
гомеостаз} — осмыслен только если уже ушли в распределённый пластичный резервуар, то есть это
инструмент \emph{той} ветки, а не нашей.
\end{intuition}

\section{Честный разбор: что покупаем и чем платим}

\paragraph{Что действительно ценно.} Ассоциативная \emph{реконструкция} вместо поиска chunk'а;
человекоподобная память (нечёткое время, ассоциативные ошибки источника, интерференция как
feature); mood как геометрия recall; CLS как защита от катастрофического забывания; и красивая
ось $e^{-tL}/\cos(t\sqrt L)/e^{-itL}$, дающая единый язык для diffusion/wave/quantum памяти.

\paragraph{Чем платим — и это надо назвать прямо.} Три технических риска и один
принципиальный.
\begin{itemize}\itemsep3pt
\item \textbf{Readout-бутылочное горлышко.} Перевод $r(t)\in\mathbb{R}^{10^5}$ обратно в
токены — самая хрупкая часть. Тупой декодер (MLP) может не извлечь смысл; умный — \emph{сам}
сделает всю работу и замаскирует бесполезность динамики (нужен экспериментальный контроль:
сначала нарочно \emph{тупой} $D$). А чтобы декодер обучить, нужны пары «состояние $\to$
правильный ответ» — которых для автобиографии \emph{нет}. Откуда берётся обучающий сигнал —
главный незакрытый вопрос. Элегантный вариант (резервуар модулирует hidden states LLM через
cross-attention/prefix) концептуально безупречен, но требует глубокого fine-tuning локальной
модели и почти исключает API-бэкенд; дешёвый вариант (декодировать в текст) — это, по сути,
RAG с лишними шагами, теряющий главную приманку «память остаётся неязыковой».
\item \textbf{Пластичность против гомеостаза (край хаоса).} Держать $\rho(W)\lesssim1$ на
\emph{непрерывно меняющемся} $W$ — нерешённая в вычислительной нейронауке задача: пластичность
воюет с устойчивостью. Плюс аттракторная память имеет малую ёмкость и \emph{ложные}
аттракторы; «плотные» версии Хопфилда чинят ёмкость, но математически превращаются в
attention — и тогда закономерен вопрос «а почему не взять память самого трансформера».
\item \textbf{Цена inference.} Прогонять $N$ шагов динамики по графу в сотни тысяч узлов на
\emph{каждый} запрос, да ещё фоном для replay — тяжело; нужна серьёзная оптимизация разреженных
умножений. Это не смертельно, но и не бесплатно (наш же принцип «граф дёшев только при чистой
обвязке»).
\item \textbf{Принципиальное: теряются свойства, ради которых мы строили граф.}
Редактируемость, забывание-по-команде (крипто-шред), провенанс, поточечная ревизия (TMS),
метки contested, семантика node\_kind — всё это распределённый резервуар \emph{уничтожает}.
Нельзя крипто-шреднуть факт из $W$; нельзя показать провенанс; нельзя гарантированно забыть
конкретное («Рин, забудь последние 10 минут» — у компаньона это \emph{UX/приватность}, а не
когнитивная поэзия). Возражение «человек тоже не забывает по команде» верно \emph{когнитивно},
но игнорирует, что у Рин есть \emph{юридическо-продуктовые} требования, которых у человека нет.
Архив-«чёрный ящик» спасает сырое удаление, но \emph{не} отменяет того, что резервуар уже
\emph{выучил} факт в $W$ — распределённо и не-разучиваемо.
\end{itemize}

\begin{principle}[Ключевая развилка: «не хранить строки» не значит «не иметь адресуемых узлов»]
Здесь легко проскочить подмену. Есть два разных смысла «памяти без фактов»:
\emph{(а)} не хранить факт как \emph{человекочитаемую строку} — узлы могут быть эмбеддингами;
это \emph{дёшево и совместимо} с нашим субстратом; \emph{(б)} вообще не иметь \emph{адресуемых
объектов памяти} — всё размазано в $W$; это и покупает человекоподобность, и \emph{ломает}
редактируемость/провенанс/забывание. Ценное — это (а) $+$ ассоциативная динамика; дорогая
переоценка — (б). А наш ретрив уже делает (а)$+$динамику: «возмути узел и смотри, что
загорится» — это диффузия по типизированному графу. Значит львиную долю пользы этой ветки мы
\emph{уже} имеем, не платя за (б). По-настоящему \emph{новых} адоптируемых кусков три:
пластичные рёбра (Hebbian-усиление от со-активации — по сути формализованная консолидация),
mood-как-нейромодулятор, и CLS-разделение быстрый/медленный (у нас грубо есть как слои
session$\to$life).
\end{principle}

\section{Квантовый слой: честно}

Соблазн прямой: ассоциативная память $=$ волновая динамика, кванты $=$ волновая механика,
значит кванты идеальны. Не совсем. \emph{Чистая} квантовая динамика унитарна
($|\psi(t)\rangle=U(t)|\psi(0)\rangle$): обратима, сохраняет норму, \emph{не забывает} и
\emph{не} схлопывается в аттрактор ($\|Ux-Uy\|=\|x-y\|$ — нет сжатия). А ассоциативной памяти
нужно ровно сжатие: много близких состояний $\to$ одно воспоминание. Поэтому квантовая
ассоциативная память неизбежно тащит \emph{диссипацию}, измерения, связь со средой,
open-system-динамику. Второй подвох: $2^n$ амплитуд $\ne$ $2^n$ \emph{доступных} чисел —
измерение выдаёт ограниченную классическую информацию, а классический вход ещё надо
\emph{загрузить} в квантовое состояние. Поэтому экспоненциальное гильбертово пространство не
даёт автоматически экспоненциальной полезной памяти.

\begin{openq}[Оценка кванта: исследовательски высоко, практически сейчас низко]
Область quantum reservoir computing и quantum associative memory реально существует и активно
развивается (спиновые сети, фотоника на сотнях мод, quantum Hopfield на малых NISQ-устройствах).
\emph{Дисклеймер о верифицируемости}: конкретные свежие эксперименты и цифры «превосходства»,
всплывшие в обсуждении, лежат за горизонтом надёжной проверки этого текста — их стоит держать
как \emph{иллюстрацию направления}, а не как установленный quantum advantage; сам advantage для
памяти пока не доказан. Разумная ближняя форма — \emph{не} «память Рин $=$ квантовый компьютер»,
а гибрид: классический гиппокамп $\to$ \emph{маленький} квантовый резервуар как чрезвычайно
богатый нелинейный feature-generator $\to$ классическое пластичное хранилище. Квант тут — не
хранилище на двадцать лет жизни, а мощный возбудитель одного оператора.
\end{openq}

\section{Вердикт и как оставить дверь открытой}

Для Рин с её требованиями (редактировать, забывать, объяснять источник, поддерживать
истинность) \textbf{символический типизированный граф остаётся хребтом}, а динамический слой —
это \emph{дополняющая} ассоциативная надстройка («что я сейчас помню») и отношения
\emph{учитель\,--\,ученик} (граф $=$ архив/учитель). Радикальную версию «убрать семантический
граф» для наших требований отклоняем — с конкретными причинами (потеря забывания/провенанса/
TMS), а не из робости. Но инженерно правильно \textbf{держать динамический слой за чистым
интерфейсом}: тогда классический решатель $R_{\text{classical}}$ можно однажды заменить на
$R_{\text{oscillator}}/R_{\text{photonic}}/R_{\text{quantum}}$, не переписывая остальное.

\begin{openq}[Фальсифицируемые эксперименты, если идти в эту ветку]
Три чистых вопроса, которые можно поставить на стенде, не ломая текущую систему:
\emph{(1)} сравнить $e^{-tL}$, $\cos(t\sqrt L)$ и $e^{-itL}$ на \emph{одной} topologии
memory-графа — какой профиль лучше поднимает правильную окрестность (diffusion vs wave vs
quantum-walk memory); \emph{(2)} даёт ли \emph{пластичная} надстройка над типизированным графом
(Hebbian-усиление со-активированных рёбер) измеримо лучший recall, чем статические веса, — и не
ценой ли катастрофической интерференции; \emph{(3)} readout-контроль: посадить нарочно тупой
декодер и проверить, есть ли вообще содержательная информация в динамике, прежде чем строить
сложный мост к LLM. Все три — на golden-микромире (гл.~\ref{ch:experiment}), где есть ground
truth.
\end{openq}
